\section{Every fixed height}\label{sec:height}

In this section the height $h$ is fixed and the width grows. We read a blackout of the $h\times m$ grid as its \emph{column word} $w=B_1B_2\cdots B_m$. Its letters lie in the alphabet $\Sigma_h$ of subsets of $\{1,\dots,h\}$. We write $|w|$ for its length and $S(w)$ for its blackout. Let $\mathcal L_h$ be the set of words with a valid blackout, the empty word included. A \emph{factor} of $w$ is a word $B_iB_{i+1}\cdots B_j$ of consecutive letters. The next proposition shows that the factors of a fixed length do not decide validity.

\begin{proposition}[validity is not a window property]\label{prop:nonlocal}
Let $g\ge1$. The blackout $S=\{(1,2),(3,2),(g+3,1),(g+3,3)\}$ of the $3\times(g+3)$ grid is axis-valid and every factor of at most $g$ columns is valid, but $S$ is not valid.
\end{proposition}

\begin{proof}
The graph $G_S$ has no cycle, and the diamond $D(1,1)$ and the rectangle on columns $2,g+3$ and rows $1,3$ both show $\{(2,1),(2,3)\}$. A factor of at most $g$ columns has at most two black points, while two rectangles with the same presentation differ in at least four black corners (Lemmas~\ref{lem:diff} and~\ref{lem:threecorners}).
\end{proof}

\begin{theorem}[regularity and rational series]\label{thm:regular}
For every $h\ge1$ the language $\mathcal L_h$ is regular, and $G_h(z,u)=\sum_{w\in\mathcal L_h}z^{|w|}u^{|S(w)|}$ is rational in $z$ and $u$.
\end{theorem}

\begin{proof}
The \emph{trace} of a rectangle $R$ is the word whose $i$th letter is the set of rows of the corners of $R$ in column $i$. It determines $C(R)$, so by Lemma~\ref{lem:diff}, $w\notin\mathcal L_h$ exactly when two different traces of length $|w|$ have letters with $Y^{(1)}_i\triangle Y^{(2)}_i\subseteq B_i$ for every $i$. The axis-aligned traces are $\emptyset^*\{r,s\}\emptyset^*\{r,s\}\emptyset^*$, and by Lemmas~\ref{lem:leftmost} and~\ref{lem:span} the tilted traces are $\emptyset^*Y_0Y_1\cdots Y_n\emptyset^*$ for finitely many templates \eqref{eq:tilt-param} of span $n\le W_h$. A finite automaton that reads $w$ and two words in parallel checks that both are traces, that they differ, and the inclusion in every column. Erasing the two words gives an automaton for the complement of $\mathcal L_h$, so $\mathcal L_h$ is regular. With the weight $u^{|B|}$ on a letter $B$, the transfer matrix $M(u)$ of a deterministic automaton for $\mathcal L_h$ gives $G_h=\alpha(I-zM(u))^{-1}\beta$, rational by Cramer's rule \cite[Proposition~I.3]{flajolet2009}.
\end{proof}

\begin{theorem}[fixed height]\label{thm:periodic}
Let $h\ge2$.
\begin{enumerate}[(a)]
\item There are integers $m_0\ge3$, $p\ge1$, and $k_0,\dots,k_{p-1}$ such that $f(h,m)=m+k_{m\bmod p}$ for every $m\ge m_0$.
\item The series $\sum_{m\ge0}a_h(m)z^m$ is rational, with the conventions of Section~\ref{sec:conv}, and the words of optimum blackouts form a regular language.

\end{enumerate}
\end{theorem}

\begin{proof}
\emph{Step 1. Bounded excess.} Let $E(w)=|S(w)|-|w|$ be the excess of $w$, and put $C_h=\max(q_h,2h-2)$. A factor $v$ of a valid word is valid, so $E(v)\le C_h$ by Theorem~\ref{thm:rowbound} when $|v|\ge3$ and by $|S(v)|\le h|v|$ otherwise. If $w=w_1w_2\in\mathcal L_h$ has $E(w)=k$, then $k-C_h\le E(w_1)\le C_h$.

\emph{Step 2. A counter.} So an automaton for $\mathcal L_h$ with a counter for $E$ that rejects outside $[k-C_h,C_h]$ recognizes the words $\mathcal L_{h,k}$ of $\mathcal L_h$ with excess $k$.

\emph{Step 3. Periodicity.} The lengths of the words of a regular language form an eventually periodic set, since the states reachable in $m$ letters evolve by a fixed map on a finite set. Take $m_0$ and $p$ that work for every $k$ from $h-1$ to $q_h$. For $m\ge3$, Lemma~\ref{lem:cross} and Theorem~\ref{thm:rowbound} make $f(h,m)-m$ the largest such $k$ with a word of length $m$, which gives (a).

\emph{Step 4. The optimum words.} So the optimum words are the words of $\mathcal L_{h,k_r}$ of length at least $m_0$ and congruent to $r$ modulo $p$, over all $r$, and finitely many shorter words. They form a regular language, with a rational series as in Theorem~\ref{thm:regular}.
\end{proof}

The theorem bounds neither $m_0$ nor $p$. On two to six rows the excess is eventually $1$, $2$, $4$, $6$, and $9$ (Theorems~\ref{thm:strip}, \ref{thm:threevalue}, \ref{thm:fourvalue}, \ref{thm:fiverow}, and~\ref{thm:sixrow}). Guilhot, Little, and Parkinson \cite[Theorem~1(3)]{glp2025} prove that the words of maximal weight in a regular language form a regular language for an additive weight bounded above. Their theorem does not give part (b) directly, since we maximize separately at each width $m$. A longer version of these proofs is the file \nolinkurl{jis/removed/sec08_height.tex} of the repository~\cite{repo}.

